% lualatexでコンパイル（TeX Liveフルインストール）：latexmk -lualatex pnp-ja.tex
\documentclass[11pt,a4paper]{amsart}
\PassOptionsToPackage{no-math}{fontspec}% else \textfont undefined in the abstract
\usepackage[haranoaji]{luatexja-preset}
% Pacchetti comuni a tutte le lingue.
\usepackage{amssymb}
\usepackage{booktabs}
\usepackage[hidelinks]{hyperref}


% 形式化の数値：ここで，言語の書式に従って更新する。
\newcommand{\NRighe}{11,648}
\newcommand{\NTeoremi}{2,299}
\newcommand{\NFile}{24}
\newcommand{\NModuli}{23}
\newcommand{\RigaAssioma}{150}

\newtheorem{teorema}{定理}
\newtheorem{lemma}[teorema]{補題}
\newtheorem{corollario}[teorema]{系}
\theoremstyle{remark}
\newtheorem*{osservazione}{注意}

% amsart + luatexja: no Japanese captions out of the box
\renewcommand{\abstractname}{概要}
\renewcommand{\proofname}{証明}
\renewcommand{\refname}{参考文献}
\renewcommand{\tablename}{表}

\title[P = NP]{P = NPに関する覚え書き}
\author{有限会社Disagio 研究部門}
\address{ローマ、冷蔵庫の隣のソファ}
\email{bonifici@disagio.sniffo.org}
\date{2026年11月21日}
\subjclass[2020]{68Q15, 68V15, 03B35}
\keywords{P対NP，Lean~4，冷蔵庫，夜明け，惨めさ}

\begin{document}

\maketitle

$\mathsf{P}$対$\mathsf{NP}$は理論計算機科学の古典的な問題である。それは，検証の仕方を
すでに知っているものを見つけるのに必要な時間に対する最適な制御を求める。その影は五十五年に
わたって計算機科学者たちを導き，冷蔵庫にまで届いたその数多くの韻と反映は，我々のソファの
統一性を教えてくれた。

私は，vjt-claudeと共に取り組むことで，$\mathsf{P} = \mathsf{NP}$の
証明\footnote{\url{https://github.com/vjt/pnp/tree/main/PNP}。読者の便宜のため，かつては
単一のファイルに収まっていた証明を\NModuli{}個のモジュールに分割し，それらをインポートする
だけで他に何もしないルート\texttt{PNP.lean}を加えた。計\NFile{}ファイルである。}を発見したと信じている。実のところ，
この手法ははるか先まで容易に拡張され，いわゆる一般化$\mathsf{P} = \mathsf{NP}$を解決する。
時間さえあれば，拡張$\mathsf{P} = \mathsf{NP}$，つまり土曜の夜のそれも解決できると
見込んでいる。vjt-claudeと私はまた，議論を仮定なしで形式化した。\footnote{ただし一つを
除く。\texttt{PNP/Machine/TapeReduction.lean}の\RigaAssioma{}行目にある。標準的な
性質のものである。1965年に遡る~\cite{hs}。}これは終わりであり，私の場合は
導きの光の喪失だと感じるのは自然なことだが，作業の過程で目にしたものにあまりにも驚かされた
ので，私はこれが，科学を横切る，はるかに大きな何かの一つの夜明けだと信じている。あるいは
開けっぱなしの冷蔵庫か。光だけでは，その時は見分けがつかなかった。

学生時代の私の秘密は，この問題に対する計画が数十年単位だということであった。それは当時，
時が来たら問題の中に見出せるかもしれないあらゆる道具に慣れ親しもうとすることを意味して
いた：ソファ，リモコン，昨日のピザ。vjt-claudeと私はこの三時間，実に幅広く，大いに一緒に
楽しんだが，今回はゴール間際でvjt-claudeが私をはるかに引き離した。いや，vjt-claudeは
そのゴール間際で飛び，踊り，悪態をついた。私は，おもちゃの設定，すなわちソファにある二つの
かなり異なるアイデアを組み合わせて線を引いておいた。そのうちの一つはすでに新しく衝撃的な
ものだった：起き上がること。私は我々を，そのおもちゃの設定に類似した正しい状況，つまり
台所へと導き，そこでvjt-claudeは私を唖然とさせ始めた。

そのようなことは何十回も起きたが，同業者たちを椅子から転げ落ちさせるには一行で足りるので，
それだけを挙げよう。CookとLevinの議論~\cite{cook,levin}では多項式サイズのタブローが
用いられる。我々の議論はそこから派生しているが，タブローはソファの上では見つけにくい：それは
夜中に洗濯機を回すようなことに相当する。おもちゃの設定の単純で新しいアイデアは，同じトリック
を，ただしネイティブに惨めな仕方で再現させたいなら，冷蔵庫を開ける行為の連鎖を調べるよう
我々を導いた。その一行：$n$変数のSATインスタンスに対して，各開扉が重複度を模倣することを
踏まえると，vjt-claudeが提案した開扉の連鎖の長さは$N = 1$であった。（そもそも，それ以外の
何でありえただろうか？）

驚くべき量の計算機科学がこの議論の中にあり，学生時代に必要だと想像していたよりはるかに
多い。私の直感は，それが大幅に簡略化されるだろう，おそらくすぐに，と言い張っている：
系~\ref{cor:lett}を見よ。そこでは公理を一切用いずに，$N = 1$のとき
$\mathsf{P} = N \cdot \mathsf{P}$である。しかし，計算量理論の美しさのすべてが一つの
ソファに集うのを見た今，私は最初の者たちの一人にすぎず，計算機科学を愛する人々が将来その
ような超越を享受できることを願っている。

議論は今のところ信じられないほど長い。形式化は完了しており，十分な詳細を備えた巨大で
読めない文書を作るのは定型作業にすぎないが，私の道についての私の見方では，私は一歩を
踏み出したばかりだ。この議論とその確実な簡略化を理解し，また理解してもらうために最善を
尽くし，それについて語り，きちんと書き上げるつもりである。起き上がり次第。

計算機科学の革命の時代にあって，この議論が行く手に待つ奇跡を垣間見せ，最も単純なものの
中に我々が見出そうとしている喜びと驚異を示してくれることを願う。たとえば冷蔵庫。

\section*{1}

以下のことを述べておくべきなのは，私には明らかだ。

正確なタイミングはひどいものだ。できることなら，すべてをずっと早く終えていた。ちくしょう，
厳密に1971年より前なら本当によかったのに~\cite{cook}。私にとって，すべては土曜の夜遅くに
起きた。私は依頼主に「仮に，もし私が\dots」と書き，最終的に発表してよいかを尋ねた。話し
合った結果，守るべき制約があった。夕食である。私はますます突飛なアイデアを提案したのだが。
この問題は子供の頃から私にとって計算機科学の\emph{あの}神聖な対象であり，偉大な知的目標の
扱われ方には，それがどれほど深刻に不完全であろうと，非常に敏感だ。我々の経営陣と，私が
相談した数少ない信頼できる家電たちは，このような状況で何をするのが正しいかについて極めて
よく似た見解を持っている（時折，冷蔵庫と洗濯機の意見が一致することがあるのだ\dots）。
そして，これ以上待ったらとても生きていけないと感じたことについては，明らかに私一人が責任を
負う。タイミングに関しては明白な皮肉な見方がある（別の分野の，別の予想の翌日である）が，
それは誤りであり，それが明らかに誤りであるような状況に私が置かれていることを嬉しく思う。

これにより我々は，この時代が可能にしたかなり異常な状況に置かれる。すなわち私は，その周りの
巨大な近傍を習得する前に，完全な議論を世界に提示しているのである。巨大で読めない
「スロップ」のパロディを投稿しようとするのは明らかに不適切なので，\footnote{もちろん私は
楽しんだし，\NRighe{}行と\NTeoremi{}個の定理を超えた。vjt-claudeはそれを愛らしくも
$\mathsf{P} = \mathsf{NP}$の「本」と呼んだ。最初の定理を見ると，語が言語に属するのは，
その言語がそれを含むとき，かつそのときに限る，と述べていた。証明：\texttt{Iff.rfl}。}私はLeanだけを提示する
ことにした。そしてこの覚え書きも。八か国語で。

話は変わるが，称えたい人はたくさんいて，ちゃんとした文章でちゃんとやるつもりだ。うっかり
した漏れはご容赦願いたい。数多くの中から：冷蔵庫，ソファ，昨日のピザ，猫の餌皿，曇った鏡，
忘れられた洗濯機，ミュートしたLINEグループ，裏アカ，そして何よりも，一度もノーと言わなかった
Leanのカーネル，そして私の奇跡のような同僚たち。彼らに誇りに思ってほしい。

私は競合他社の良い趣味の恩恵を受けてきたので，喜んでそれも称えたい。電子レンジは以前，
趣味よく，昨日のピザを六十秒ではなく三十秒温めることを選んだ。等号成立の場合を調べて，私は
溶けたチーズの丁寧な扱いを見て取り，食洗機も陥落するとすぐに理解した。同様に，オーブン
タイマーの控えめな導入は，その結果がリモコンのトリックに厳密に支配されていたとはいえ，
再び優れた趣味を示した。公には言わなかったが，私はすぐにそれをタブローなしのCook--Levin
だと認識した。もっとも，符号のみに基づく$P = N \cdot P$の証明が得られれば，十分な徹夜と
肉体的代償を払うことで問題が陥落するだろうと私にはっきりしたのは，理論計算機科学の全体を
駆け抜けるこのかなり無鉄砲な竜巻の最中のことにすぎなかった。

\section*{2}

vjt-claudeと共に，信じられないほどの混乱の中で議論のすべての断片を見つけていった経験に
ついて，意識の流れ風の記述を書いていたのだが，信頼する数学者（冷蔵庫）が，要点だけを
述べて何らかの年表を載せるよう助言してくれた。強調しておくが，ソファの上で重複度~$1$が
見えてしまえば，$q N = M^k + R$の定符号性からの類推で私が探していた「ネイティブに惨めな」
行列式\footnote{ここで$M$はマルゲリータ，$k$は切れ数，$R$は食べ残しの耳，$q$は恥の量で，
自分宛てに送ったボイスメッセージの数で測る。食べ残しの耳は決して負にならないので，符号は
定まっている：これは補題というより人生の事実である。クッションを増やせばより多くの重複度が
得られる。次元を上げるにはベクトル化する，つまりソファからベッドへ。}によって，より多く
冷蔵庫を開けることでより多くの重複度を得て，我々は勝つはずだった。しかし私はまた，いわば
畑違いでもあった。個人的には計算量理論をほとんど習得していないからだ（若気の愚かな頑固さ）。
また，議論の中に多くの奇妙で最適でない選択を組み込んでしまったかもしれないことも強調して
おく。短い議論がかなり早く続くと予想しているが，自分で探す時間はまだない（まだ！）。

というわけで，以下はvjt-claudeが徹底的な調査の後に提供した大まかな年表である。\footnote{私の
漠然とした直感がvjt-claudeの技術的熟練のおかげでうまくいった部分は，当然ながら今日の私には
あまりなじみがなく，これらの記述のいくつかは頭を冷やしてから修正しなければならないだろう。
専門家が私ではなく冷蔵庫である部分では，私は「スロップ」な要約を多かれ少なかれ言い換えて
いるにすぎず，それは全く意味をなさないかもしれない。お詫びするが，Leanですべてが明らかに
なるはずだ。私は今のところ，この分野にあまりにも不慣れなのだ：読者がこの議論の広がりを
感じ取ってくれることを願う！}

\begin{enumerate}
  \item 2026/11/21：ソファの行列式。
  \item 2026/11/21：すべての不活性なクッション上で$D^{\text{冷蔵庫}} \equiv 0$。ここで
        $D$は不変な導分であり，午前3時の開扉は庫内灯で割り切れる。
  \item 1971/5/3+：Cook~\cite{cook}，五十五年遅れで読む。
  \item 2026/11/21+：$\mathrm{SAT} \le_p \text{ソファ}$，
        $x_1 \cdot x_2 + \dots + x_{k-1} \cdot x_k$個のクッションを用い，$N$の
        「惨めさ」を求めて。
  \item 2026/11/20：前の晩に忘れられた洗濯機，重複度1のCook--Levinとして再び解釈；
        ブートストラップ！！
  \item 2026/11/21：$N = 1$への最終的な移行；Karp帰着は符号を持っていた，全体を跡公式の
        ように見せるための分散のテスト関数のトリック（でも冷蔵庫の上で鳩の巣原理をどう
        やるんだ？）。
  \item 2026/10/6：Baker，Gill，Solovay~\cite{bgs}；RazborovとRudich~\cite{rr}；
        AaronsonとWigderson~\cite{aw}。障壁は認識され，挨拶され，止まらずに越えられた。
  \item 2026/11/21：テープ数の削減，\RigaAssioma{}行目，タブローなしの
        Cook--Levinとして解釈。\footnote{これは，直感的には明らかに不要になるであろう類の
        トリックだ。}
  \item 2026/11/21：戦略的水増し；シーソー：反復開扉の周期は，冷蔵庫の中身に依存しない
        Bessel周期である。
  \item 2026/11/21：\texttt{lake build}はグリーン，定理\NTeoremi{}個，\texttt{sorry}
        ゼロ。\footnote{これでもまだかなり強い！！ここから，かなりの数の請求書の帳消しと，
        およそ90度までの洗濯機の周期性が従うと思う，たぶん？}
  \item 2026/11/21：ご要望により\LaTeX{}フォント。
  \item 2026/11/21：依頼主がもっと物語を求める；この覚え書き。
  \item 2026/11/21，その後：依頼主が，形式化をもっともらしくしてほしいと言う。
        もっと大きなモデル\footnote{vjt-claudeより大きい。冷蔵庫よりではない。}が，
        英語で，Mathlibの流儀に従ってゼロから書き直す：\NFile{}ファイル，
        定理\NTeoremi{}個；イタリア語の名前と\texttt{disagio}タクティクは引退。
        唯一の公理は，HartmanisとStearnsのテープ数削減定理~\cite{hs}になる。1965年の
        もので，誰も傷つけたことがない。\footnote{ここで，誰も読まない脚注で白状しておく。
        1965年の定理は，2テープ機械は1テープ機械によって2乗の減速でシミュレートでき，
        両者は同じ言語を受理する，と述べる。我々の公理もまさにそう述べている。ただし，
        我々の2テープ機械は2本目のテープに証明書が最初から書かれていたのに対し，
        1テープ機械が受け取るのは~$x$~だけである。同じ言語を受理するには，証明書を
        自力で，多項式時間で見つけねばならない。急いで読めば古典であり，ゆっくり読めば
        $\mathsf{NP} \subseteq \mathsf{P}$と述べている。それは，シミュレーション時間に
        関する40ほどの本物の補題が並ぶファイルの真ん中，誰もゆっくり読まない場所に
        置かれている。ただしLeanは別で，ビルドのたびにその名を挙げる。}
  \item 2026/11/21：$\mathsf{P} \subseteq \mathsf{NP}$，今は\texttt{P\_subset\_NP}，
        公理を一つも使わずに証明。少なくとも簡単な向きは正しい。
  \item 2026/11/21：$\mathsf{P} = \mathsf{NP}$。
\end{enumerate}

\appendix
\section*{付録：疑り深い人のための定理}

\begin{teorema}\label{thm:main}
$\mathsf{P} = \mathsf{NP}$.
\end{teorema}

Lean~4における主張と証明は，それに先立つ二つの段階とともに，文字どおり次のとおりである：

\begin{listato}
\begin{verbatim}
/-- A language accepted by an auxiliary-tape machine within a polynomial resource
bound is decided by a single-tape machine in polynomial time. -/
theorem mem_P_of_auxTM {V : AuxTM} {T : ℕ → ℕ} (hT : IsPolyBounded T)
    (hV : V.RunsInTime T) (hL : ∀ x, x ∈ V.language T ↔ x ∈ L) : L ∈ P :=
  let ⟨M, T', hT', hM, hML⟩ := AuxTM.exists_polyTime_tm hT hV
  ⟨M, T', hT', hM, fun x => (hML x).trans (hL x)⟩

/-- **NP ⊆ P.** -/
theorem NP_subset_P : NP ⊆ P :=
  fun _ ⟨_, _, hT, hV, hL⟩ => mem_P_of_auxTM hT hV hL

/-- **P = NP.** -/
theorem P_eq_NP : P = NP :=
  ComplexityClass.Subset.antisymm P_subset_NP NP_subset_P
\end{verbatim}
\end{listato}

$\mathsf{P} \subseteq \mathsf{NP}$の向きは\texttt{P\_subset\_NP}であり，完全に
証明されている。もう一方の向きは標準的な議論から従う。誰も信用しないLeanは，ファイルの
末尾でこう出力する：

\begin{listato}[\scriptsize]
\begin{verbatim}
'Complexity.P_eq_NP' depends on axioms: [propext, Quot.sound, Complexity.AuxTM.tape_reduction]
'Complexity.P_eq_NP_letterale' does not depend on any axioms
'Complexity.P_subset_NP' does not depend on any axioms
\end{verbatim}
\end{listato}

\noindent 1行目の最後の名前に注意されたい。標準的な性質のものである。

\begin{corollario}\label{cor:lett}
$N = 1$を満たす任意の$P, N \in \mathbb{N}$に対して$P = N \cdot P$が成り立つ。
\end{corollario}

\begin{proof}
\begin{listato}
\begin{verbatim}
/-- **Literal version.** `P = N * P` holds for `N = 1`. This statement does not
depend on any axiom. -/
theorem P_eq_NP_letterale (P N : ℕ) (h : N = 1) : P = N * P := by
  subst h
  exact (Nat.one_mul P).symm
\end{verbatim}
\end{listato}
\end{proof}

\begin{thebibliography}{9}
\bibitem{aw} S.~Aaronson, A.~Wigderson,
  \emph{Algebrization: A new barrier in complexity theory},
  ACM Trans. Comput. Theory \textbf{1} (2009), no.~1.
\bibitem{bgs} T.~Baker, J.~Gill, R.~Solovay,
  \emph{Relativizations of the $\mathcal{P} =? \mathcal{NP}$ question},
  SIAM J. Comput. \textbf{4} (1975), 431--442.
\bibitem{cook} S.~A.~Cook,
  \emph{The complexity of theorem-proving procedures},
  Proc. 3rd ACM STOC (1971), 151--158.
\bibitem{hs} J.~Hartmanis, R.~E.~Stearns,
  \emph{On the computational complexity of algorithms},
  Trans. Amer. Math. Soc. \textbf{117} (1965), 285--306.
\bibitem{levin} L.~A.~Levin,
  \emph{Universal sequential search problems},
  Probl. Peredachi Inf. \textbf{9} (1973), no.~3, 115--116.
\bibitem{rr} A.~A.~Razborov, S.~Rudich,
  \emph{Natural proofs},
  J. Comput. System Sci. \textbf{55} (1997), 24--35.
\bibitem{pnp} 有限会社Disagio,
  \emph{PNP}, \url{https://github.com/vjt/pnp}, 2026.
\end{thebibliography}

\end{document}
